\paragraph{Hint for Exercise 1: Universal triangle}\mbox{}\par
Write the commuting triangle and emphasize that the target must already be a sheaf.


\paragraph{Hint for Exercise 2: Uniqueness of associated sheaf}\mbox{}\par
Factor each unit through the other sheafification, then use uniqueness twice.


\paragraph{Hint for Exercise 3: Basis for the espace etale}\mbox{}\par
Every germ has a representative.  For the intersection condition, use equality of germs to shrink to a common neighborhood.


\paragraph{Hint for Exercise 4: Local homeomorphism}\mbox{}\par
The inverse is $x\mapsto s_x$; use the definition of the basis topology.


\paragraph{Hint for Exercise 5: Local representability}\mbox{}\par
Continuity into a basis neighborhood gives a local representative; conversely local continuity is enough.


\paragraph{Hint for Exercise 6: Associated sections form a sheaf}\mbox{}\par
Glue the underlying maps $U_i\to E_{\mathcal F}$ and use that continuity is local.


\paragraph{Hint for Exercise 7: The canonical morphism}\mbox{}\par
Compare the germ of $s|_V$ at $x$ with the germ of $s$ at $x$.


\paragraph{Hint for Exercise 8: Surjectivity on stalks}\mbox{}\par
Represent an arbitrary element of $\mathcal F_x$ by $s\in\mathcal F(U)$.


\paragraph{Hint for Exercise 9: Injectivity on stalks}\mbox{}\par
Evaluate a germ of a sheafified section at $x$, then use local representability and equality of original germs.


\paragraph{Hint for Exercise 10: Factorization existence}\mbox{}\par
Choose a cover on which the sheafified section comes from $s_i\in\mathcal F(U_i)$ and glue $u(s_i)$.


\paragraph{Hint for Exercise 11: Factorization uniqueness}\mbox{}\par
Two factors agree on every locally representable neighborhood.


\paragraph{Hint for Exercise 12: Functoriality}\mbox{}\par
Apply the universal property to $\eta_\mathcal G\circ f$ and use uniqueness for identities and composites.


\paragraph{Hint for Exercise 13: Naturality of the unit}\mbox{}\par
Both sides are maps from $\mathcal F$ to the sheaf $\mathcal G^\#$ with the same defining factorization.


\paragraph{Hint for Exercise 14: Stalk map of a sheafified morphism}\mbox{}\par
Take stalks of the naturality square and cancel the two unit isomorphisms.


\paragraph{Hint for Exercise 15: Sheafifying a sheaf}\mbox{}\par
Factor $\id_\mathcal G$ through $\eta_\mathcal G$; then use uniqueness to get a two-sided inverse.


\paragraph{Hint for Exercise 16: Idempotence}\mbox{}\par
Apply the preceding exercise to the sheaf $\mathcal F^\#$.


\paragraph{Hint for Exercise 17: Constant presheaf on two points}\mbox{}\par
The two stalks are $A$; use the classification of sheaves on a discrete space.


\paragraph{Hint for Exercise 18: Bounded continuous functions}\mbox{}\par
Every continuous function is locally bounded.  Use local representability for both inclusions.


\paragraph{Hint for Exercise 19: A unit that is not injective}\mbox{}\par
Every germ is zero because every point has a proper neighborhood on which the section restricts to zero.


\paragraph{Hint for Exercise 20: Restriction to an open}\mbox{}\par
Compare the two espace etale spaces over $U$ or use the universal property after restriction.


\paragraph{Hint for Exercise 21: Cofinal basis neighborhoods}\mbox{}\par
Basis neighborhoods are cofinal among all neighborhoods of $x$.


\paragraph{Hint for Exercise 22: Separated presheaf}\mbox{}\par
Separated means compatible equality on a cover determines the global section uniquely.


\paragraph{Hint for Exercise 23: First plus is separated}\mbox{}\par
Represent two plus-sections by compatible families; refine covers until their local representatives agree.


\paragraph{Hint for Exercise 24: Plus of a separated presheaf}\mbox{}\par
A compatible family of plus-sections can be represented on refinements by original sections; separatedness makes all choices compatible.


\paragraph{Hint for Exercise 25: Two pluses}\mbox{}\par
Apply the preceding two facts first to $\mathcal F$ and then to $\mathcal F^+$.  Verify the same universal mapping property.


\paragraph{Hint for Exercise 26: Adjunction}\mbox{}\par
Given a map from $\mathcal F$ to a sheaf, send it to its unique factor through $\eta$; the inverse is precomposition with $\eta$.


